\section{Conclusión} \label{sec:conclusion}

Este trabajo estima el efecto de la apertura de una estación de servicio sobre el precio de las incumbentes de su mercado local en Chile, usando los precios históricos de bencinaenlinea.cl entre 2012 y 2026 y un estudio de eventos escalonado que se estima tanto por efectos fijos bidireccionales como con el estimador de \textcite{SUNABRAHAM2021}. La identificación descansa en que, controlando por efectos fijos de estación, mes, región por año y marca por año, el momento en que abre una estación no responde a la evolución de los precios de las incumbentes \parencite{ARCIDIACONO2020,FISCHER2025}. Las tres hipótesis planteadas encuentran respaldo, con matices que dependen del combustible.

En primer lugar, la entrada de una competidora a 2 kilómetros o menos reduce el precio de las incumbentes. El efecto promedio es de 0,40\% en las gasolinas 93 y 95, de 0,24\% en la gasolina 97 y de 0,59\% en el diésel. En las gasolinas 93 y 95 y en el diésel la caída no se revierte en los años siguientes a la entrada, mientras que en la gasolina 97 deja de distinguirse de cero desde el tercer semestre. En segundo lugar, el efecto decrece con la distancia a la entrante, es significativo hasta los 3 kilómetros en los cuatro combustibles y no se distingue de cero entre los 3 y los 5 kilómetros. En tercer lugar, el ajuste no es homogéneo dentro del mercado. El percentil 10 de los precios del mercado local cae más que el percentil 90, 0,62\% frente a 0,29\% en la gasolina 93, 0,57\% frente a 0,24\% en la 95 y 0,64\% frente a 0,46\% en el diésel, y la regresión de distribución muestra que el desplazamiento se concentra en la mitad baja de la distribución. La entrada, por tanto, aumenta la dispersión de precios del mercado local. La gasolina 97, el combustible de menor volumen, es la excepción, pues ningún cuantil del mercado responde de forma significativa y el desplazamiento de la distribución se ubica en su centro.

Estos resultados se sostienen en los ejercicios de robustez. El efecto del primer semestre sobre el precio de las incumbentes resiste una relajación de tendencias paralelas tan grandes como el mayor salto observado antes de la entrada en los cuatro combustibles \parencite{RAMBACHAN2023}. El número de incumbentes no cambia durante el primer año, cuando ya se observa la caída de precios, por lo que el efecto no proviene de la salida de estaciones. Al adelantar seis meses la fecha de entrada el quiebre sigue ocurriendo en la fecha registrada, y las entradas asignadas al azar a estaciones de control generan efectos cercanos a cero, con valores $p$ por aleatorización del efecto real de 0,002 o menos en las gasolinas 93 y 95 y en el diésel.

La evidencia tiene tres implicancias para la práctica de la autoridad de competencia. La primera se refiere a la delimitación del mercado geográfico. El umbral estimado respalda el radio de 3 kilómetros que la \textcite{FNE2021informe} ha considerado adecuado. Sin embargo, en las incumbentes ubicadas junto a autopistas la entrada no reduce el precio, lo que sugiere que para ellas un radio fijo no es el criterio apropiado y que su mercado debiera definirse a lo largo del corredor vial. La segunda se refiere a las barreras a la entrada. Si la escasez de terrenos aptos y las exigencias de los planos reguladores documentadas por el \textcite{TDLC2012} y la \textcite{FNE2021informe} impiden una apertura, las incumbentes a menos de 3 kilómetros mantienen precios mayores de forma persistente, de modo que la regulación de uso de suelo tiene consecuencias de competencia que debieran considerarse en su diseño. La tercera se refiere a la transparencia de precios. \textcite{LUCO2019} muestra que la publicación obligatoria de precios en línea elevó los márgenes en este mismo mercado. Los resultados de este trabajo indican que, en presencia de esa transparencia, la entrada reduce sobre todo los precios de las estaciones más baratas. Esto es consistente con que la ganancia de la entrada sea mayor para los consumidores que comparan precios, en línea con lo que \textcite{FISCHER2025} encuentran para Alemania.

El estudio tiene las siguientes limitaciones. El supuesto de exogeneidad del momento de la entrada no puede verificarse directamente, y aunque los coeficientes previos a la entrada son pequeños en la especificación principal, en el diésel los precios del mercado local ya venían cayendo antes de la entrada por lo que parte del efecto estimado en esos casos podría reflejar una tendencia previa. El efecto posterior a la pandemia es menor y menos persistente, y este trabajo no logra identifica si ello responde a la madurez del mercado, a la menor frecuencia de actualización del MEPCO desde 2023 o a una composición distinta de las entradas. Por último, la base no contiene volúmenes de venta, lo que impide traducir las caídas de precio en ganancias de bienestar para los consumidores o identificar qué consumidores compran en las estaciones cuyos precios caen más.

